\begin{HVTPage}{Capítulo 2}{Electrólisis}{Ruta electroquímica}
\HVTLead{La electrólisis separa el agua en hidrógeno y oxígeno mediante electricidad. El diseño debe considerar calidad del agua, potencia, eficiencia, temperatura, presión y tecnología del electrolizador.}
\begin{HVTFlow}\HVTFlowItem{1}{Agua tratada}{Filtración, desmineralización y control de conductividad.}\HVTFlowItem{2}{Electrolizador}{Conversión electroquímica y separación inicial.}\HVTFlowItem{3}{Secado}{Eliminación de humedad, medición y control de pureza.}\end{HVTFlow}
\begin{HVTTable}{Variable}{Símbolo}{Unidad}
\HVTTableRow{Corriente eléctrica}{$I$}{A}
\HVTTableRow{Eficiencia faradaica}{$\eta_F$}{--}
\HVTTableRow{Constante de Faraday}{$F$}{C/mol}
\HVTTableRow{Potencia eléctrica}{$P$}{kW}
\end{HVTTable}
\HVTEquation{\dot{n}_{H_2}=\frac{\eta_F\cdot I}{2\cdot F}}{}
\HVTText{La expresión se deriva de la ley de Faraday para una reacción que requiere dos electrones por mol de H2. Se adopta F = 96 485,332 12 C/mol, valor exacto reportado por NIST. \HVTCite{1} La revisión de Carmo et al. documenta los fundamentos y variables de la electrólisis PEM. \HVTCite{2}}
\end{HVTPage}

\begin{HVTPage}{Capítulo 2}{Ejemplo E-1}{Estimación con corriente conocida}
\begin{HVTExample}{Ejemplo E-1}{Estimación con corriente conocida}\begin{enumerate}\item Definir la corriente estable del módulo. \item Establecer la eficiencia faradaica esperada. \item Calcular el flujo molar de H2. \item Convertir el resultado a caudal volumétrico o másico. \item Verificar capacidad de secado, venteo y almacenamiento.\end{enumerate}\end{HVTExample}
\end{HVTPage}
